% GENERATED FILE. Edit canonical lessons, then run npm run generate:books.
% canonical-source-hash: fnv1a64:d907cac91db0430e
% canonical-lessons: ZH-C56-wan, ZH-C56-daqiu, ZH-C56-yundong, ZH-C56-paobu, ZH-C56-xiao

\chapter{Free Time}
\label{ch:zh-free-time}
\hlchaptermodality{56}

\begin{chapteropening}
\textbf{By the end of this chapter:} \emph{I can say what I do for fun: play, sport, running and laughing.}

Complete the last lesson of chapter 56: I can say what I do for fun: play, sport, running and laughing.
\end{chapteropening}

\section[wán]{wán — to play; to have fun}

\label{lesson:ZH-C56-wan}



\begin{quote}
Before the new one: say the Mandarin for to sing, then the Mandarin for coffee.
\end{quote}

\subsection*{You'll want to know: wán}
\emph{wán} — \textquotedblleft{}to play; to have fun\textquotedblright{}.

\begin{quote}
\emph{háizi xǐhuan wán} — Children love to play.
\end{quote}

Its characters come later; for now it is a word for the ear.

\begin{cousinweb}
\emph{wán} is inherited from Middle Chinese \emph{ngwanH} (Guangyun, in Baxter's transcription), glossed 'to toy with'.
\end{cousinweb}

\begin{grammarlens}[title={What you've built}]
One more word for everyday answers.
\end{grammarlens}

\subsection*{Your turn}
\begin{itemize}
\raggedright
  \item \emph{Say it:} \emph{wán}
  \item \emph{Say it:} \emph{wán}, once more
  \item \emph{Say it:} \emph{háizi xǐhuan wán} — Children love to play.
  \item \emph{Recall:} say the Mandarin for to sing, then the Mandarin for coffee, then say \emph{wán} again
\end{itemize}

\begin{tcolorbox}[breakable,colback=teal!4,colframe=teal!35!black,title={Before you move on}]
What does wán mean? (\textquotedblleft{}To play; to have fun\textquotedblright{}.) Say it once more.
\end{tcolorbox}
\section[dǎqiú]{dǎqiú — to play ball}

\label{lesson:ZH-C56-daqiu}



\begin{quote}
Before the new one: say the Mandarin for coffee, then the Mandarin for to play.
\end{quote}

\subsection*{You'll want to know: dǎqiú}
\emph{dǎqiú} — \textquotedblleft{}to play ball\textquotedblright{}.

\begin{quote}
\emph{xīngqītiān wǒmen dǎqiú} — We play ball on Sunday.
\end{quote}

Its characters come later; for now it is a word for the ear.

\begin{cousinweb}
\emph{dǎqiú} is built from dǎ 'to strike' and qiú 'ball', both inherited from Middle Chinese, \emph{taengX} and \emph{gjuw} (Guangyun, in Baxter's transcription), though dǎ is not the sound \emph{taengX} would regularly give. The dictionary glosses \emph{gjuw} as 'fine jade', and under another way of writing it as a ball of fur that is struck.
\end{cousinweb}

\begin{grammarlens}[title={What you've built}]
One more word for everyday answers.
\end{grammarlens}

\subsection*{Your turn}
\begin{itemize}
\raggedright
  \item \emph{Say it:} \emph{dǎqiú}
  \item \emph{Say it:} \emph{dǎqiú}, once more
  \item \emph{Say it:} \emph{xīngqītiān wǒmen dǎqiú} — We play ball on Sunday.
  \item \emph{Recall:} say the Mandarin for coffee, then the Mandarin for to play, then say \emph{dǎqiú} again
\end{itemize}

\begin{tcolorbox}[breakable,colback=teal!4,colframe=teal!35!black,title={Before you move on}]
What does dǎqiú mean? (\textquotedblleft{}To play ball\textquotedblright{}.) Say it once more.
\end{tcolorbox}
\section[yùndòng]{yùndòng — sport; exercise}

\label{lesson:ZH-C56-yundong}



\begin{quote}
Before the new one: say the Mandarin for to play, then the Mandarin for to play ball.
\end{quote}

\subsection*{You'll want to know: yùndòng}
\emph{yùndòng} — \textquotedblleft{}sport; exercise\textquotedblright{}.

\begin{quote}
\emph{wǒ xǐhuan yùndòng} — I like sport.
\end{quote}

Its characters come later; for now it is a word for the ear.

\begin{cousinweb}
\emph{yùndòng} is built from yùn 'to move round, to carry' and dòng 'to move', both inherited from Middle Chinese, \emph{hjunH} and \emph{duwngX} (Guangyun, in Baxter's transcription). The sense 'sport' is not given there.
\end{cousinweb}

\begin{grammarlens}[title={What you've built}]
One more word for everyday answers.
\end{grammarlens}

\subsection*{Your turn}
\begin{itemize}
\raggedright
  \item \emph{Say it:} \emph{yùndòng}
  \item \emph{Say it:} \emph{yùndòng}, once more
  \item \emph{Say it:} \emph{wǒ xǐhuan yùndòng} — I like sport.
  \item \emph{Recall:} say the Mandarin for to play, then the Mandarin for to play ball, then say \emph{yùndòng} again
\end{itemize}

\begin{tcolorbox}[breakable,colback=teal!4,colframe=teal!35!black,title={Before you move on}]
What does yùndòng mean? (\textquotedblleft{}Sport; exercise\textquotedblright{}.) Say it once more.
\end{tcolorbox}
\section[pǎobù]{pǎobù — to run; to jog}

\label{lesson:ZH-C56-paobu}



\begin{quote}
Before the new one: say the Mandarin for to play ball, then the Mandarin for sport.
\end{quote}

\subsection*{You'll want to know: pǎobù}
\emph{pǎobù} — \textquotedblleft{}to run; to jog\textquotedblright{}.

\begin{quote}
\emph{wǒ zǎoshang pǎobù} — I go for a run in the mornings.
\end{quote}

Its characters come later; for now it is a word for the ear.

\begin{cousinweb}
\emph{pǎobù} is built from pǎo 'to run' and bù 'step', both inherited from Middle Chinese, \emph{baew} and \emph{buH} (Guangyun, in Baxter's transcription); the dictionary glosses \emph{baew} as a foot pawing the ground.
\end{cousinweb}

\begin{grammarlens}[title={What you've built}]
One more word for everyday answers.
\end{grammarlens}

\subsection*{Your turn}
\begin{itemize}
\raggedright
  \item \emph{Say it:} \emph{pǎobù}
  \item \emph{Say it:} \emph{pǎobù}, once more
  \item \emph{Say it:} \emph{wǒ zǎoshang pǎobù} — I go for a run in the mornings.
  \item \emph{Recall:} say the Mandarin for to play ball, then the Mandarin for sport, then say \emph{pǎobù} again
\end{itemize}

\begin{tcolorbox}[breakable,colback=teal!4,colframe=teal!35!black,title={Before you move on}]
What does pǎobù mean? (\textquotedblleft{}To run; to jog\textquotedblright{}.) Say it once more.
\end{tcolorbox}
\section[xiào]{xiào — to laugh; to smile}

\label{lesson:ZH-C56-xiao}



\begin{quote}
Before the new one: say the Mandarin for sport, then the Mandarin for to run.
\end{quote}

\subsection*{You'll want to know: xiào}
\emph{xiào} — \textquotedblleft{}to laugh; to smile\textquotedblright{}.

\begin{quote}
\emph{tā xiào le} — She laughed.
\end{quote}

Its characters come later; for now it is a word for the ear.

\begin{cousinweb}
\emph{xiào} is inherited from Middle Chinese \emph{sjewH} (Guangyun, in Baxter's transcription), glossed 'glad, joyful'.
\end{cousinweb}

\begin{grammarlens}[title={What you've built}]
One more word for everyday answers.
\end{grammarlens}

\subsection*{Your turn}
\begin{itemize}
\raggedright
  \item \emph{Say it:} \emph{xiào}
  \item \emph{Say it:} \emph{xiào}, once more
  \item \emph{Say it:} \emph{tā xiào le} — She laughed.
  \item \emph{Recall:} say the Mandarin for sport, then the Mandarin for to run, then say \emph{xiào} again
\end{itemize}

\begin{tcolorbox}[breakable,colback=teal!4,colframe=teal!35!black,title={Before you move on}]
What does xiào mean? (\textquotedblleft{}To laugh; to smile\textquotedblright{}.) Say it once more.
\end{tcolorbox}
